\section{Conclusion}\label{sec:conclusion}

The central distinction is between
\[
\text{same endpoint judgement} \;\neq\; \text{same warrant}.
\]
Grounded transport carries evidence; composition and contexts preserve it, at the level of crossings always and at the level of complete certificates when an evidence lift exists; erasure yields ordinary substitutability and loses information irreversibly; observational agreement cannot license lineage-sensitive substitution, and located obstructions explain failed transport; open assessments preserve unresolved obligations, and a refutation does not discard them; a finite lineage audit is proved to reflect its specification and, once extracted, issues certificates; and extraction retains the evidential data that downstream audits require.

Ordinary rewriting answers whether one expression may replace another under a relation. Grounded transport additionally answers through which process, under which lineage, within what coverage, and with which remaining obligations that replacement is warranted. Erasure recovers ordinary substitutability, but ordinary substitutability cannot reconstruct the warrant that was erased.

The calculus does not make declarations true. It makes precisely explicit what has been declared, checks what can be checked, and reports what has not been resolved. Witness-carrying rewriting tactics, and a treatment of certificates carried along evolutions, are natural next steps that rest on the semantics established here.
